\Needspace{14\baselineskip}
\fancyhead[L]{\small SECOND COMET OF 1805.}
\begin{center}
SECOND COMET OF THE YEAR 1805.

\rule{0.6\linewidth}{0.4pt}

Monatliche Correspondenz. Vol. XIII. pp. 83–91. January 1806.

\rule{0.6\linewidth}{0.4pt}

—\quad—\quad—
\end{center}

Dr. \textsc{Gauss} observed this comet on 8 December. It was compared with 359 in Aquarius and another star from the \emph{Hist. célest.}

\begin{center}
\begingroup\fontsize{9}{11}\selectfont\setlength{\tabcolsep}{3pt}\renewcommand{\arraystretch}{1.1}
\begin{tabular}{r|r|r|r}
1805 & \shortstack{Mean time in\\Brunswick} & \shortstack{Apparent right\\ascension} & \shortstack{Apparent southern\\declination}\\ \hline
8 December & $6^{\mathrm u}\ 55^{\mathrm m}\ 17^{\mathrm s}$ & $353^\circ\ 7'\ 40''$ & $23^\circ\ 36'\ 24''$ \\
8 & $7\ 58\ 3$ & $352\ 57\ 57$ & $24\ 6\ .\ .$ \\
\end{tabular}\endgroup
\end{center}

\par\indent The second declination is only an estimate. From this observation, from \textsc{Bouvard}’s first observation of 16 November, and from Dr. \textsc{Olbers}’s observation of 2 December given above, Dr. \textsc{Gauss} calculated the following parabolic elements.

\begin{center}
\begingroup\fontsize{8}{10}\selectfont\setlength{\tabcolsep}{3pt}\renewcommand{\arraystretch}{1.1}
\begin{tabular}{lr}
Passage through perihelion on 31 December 1805 , & $7^{\mathrm u}\ 20^{\mathrm m}\ 39^{\mathrm s}$ Seeberg time \\
Longitude of perihelion , . . . . . . . . . . . . & $109^\circ\ 23'\ 40''$ \\
Longitude of the ascending node . . . . . . . . . . & $250\ 33\ 14$ \\
Inclination of the orbit . , . . . . . \raisebox{.3ex}{.} . . . . . . . & $16\ 33\ 33$ \\
Logarithm of the least distance . . . . . . . . & $9.9502477$ \\
Motion . . . . . . . . . . . . . . . . . . & prograde. \\
\end{tabular}\endgroup
\end{center}

\Needspace{24\baselineskip}
\par\indent Following the method in the article in the May 1804 issue of the M. C. [pp. \hyperlink{native-99}{94} and \hyperlink{native-125}{120} of this volume], the coordinates are given by the following formulas.

\[
\begin{gathered}
x = \frac{\alpha \sin(v + 198^\circ\,37'\,20'')}{\cos\frac{1}{2}v^2},\\[4pt]
y = \frac{\ObservedBeta \sin(v + 103^\circ\,41'\,12'')}{\cos\frac{1}{2}v^2},\\[4pt]
z = \frac{\gamma \sin(v + 148^\circ\,41'\,54'')}{\cos\frac{1}{2}v^2}
\end{gathered}
\]

\noindent where $v$ denotes the anomaly of the comet, and Log. $\alpha = 9.933968$, Log. $\ObservedBeta = 9.930525$, Log. $\gamma = 9.551021$.

Using his elements and allowing for parallax, Dr. \textsc{Gauss} finds the comet’s distance from Earth to be 0.24990 on 16 Nov., 0.08913 on 2 Dec., 0.04850 at the first observation on 8 Dec., and 0.04841 at the second observation on 8 Dec.

The brightness, on the other hand, is 16.6 on 16 November, 117.8 on 2 December, and 439.2 on 8 December, if the brightness at distance 1 from Earth and from the Sun is taken as the unit.

\begin{center}
—\quad—\quad—
\end{center}

\hfill[\textsc{von Zach.}]
